\documentclass[11pt]{article}
\usepackage[a4paper,margin=2.5cm]{geometry}
\usepackage{amsmath}
\usepackage{amssymb}
\usepackage{enumitem}
\usepackage{parskip}
\setlist[enumerate,2]{label=(\alph*), itemsep=1pt, topsep=2pt}
\setlength{\parindent}{0pt}
\newcommand{\correct}[1]{\textbf{#1}\;$\checkmark$}
\begin{document}
\section*{Subgame perfection}
\begin{enumerate}[itemsep=6pt]
  \item A subgame perfect equilibrium is a Nash equilibrium that ...
  \begin{enumerate}
    \item induces a Nash equilibrium in every subgame, including unreached ones
    \item reaches the last node of the tree
    \item maximises the sum of the players' payoffs
    \item uses only pure strategies
  \end{enumerate}
  \item Backward induction solves a perfect-information game by ...
  \begin{enumerate}
    \item converting it to a zero sum game
    \item working from the final decisions back to the root, choosing optimally at each node
    \item trying every strategy profile at random
    \item ignoring the players' payoffs
  \end{enumerate}
  \item A non-credible threat is one that ...
  \begin{enumerate}
    \item is written down before the game begins
    \item the threatening player would not actually want to carry out if the subgame were reached
    \item both players agree to
    \item always leads to the highest payoff
  \end{enumerate}
  \item Which is true of Nash equilibria versus subgame perfect equilibria?
  \begin{enumerate}
    \item subgame perfect equilibria need not be Nash equilibria
    \item every Nash equilibrium is subgame perfect
    \item they are always the same
    \item every subgame perfect equilibrium is a Nash equilibrium, but not conversely
  \end{enumerate}
  \item Every finite game of perfect information has ...
  \begin{enumerate}
    \item exactly two equilibria
    \item no Nash equilibrium
    \item a Nash equilibrium in pure strategies, found by backward induction
    \item a unique mixed equilibrium only
  \end{enumerate}
  \item In the centipede game, backward induction predicts that ...
  \begin{enumerate}
    \item both players pass to the final leaf
    \item there is no subgame perfect equilibrium
    \item the first player takes immediately, even though both could do better by passing
    \item the players alternate forever
  \end{enumerate}
  \item A node initiates a subgame when ...
  \begin{enumerate}
    \item both players move there at once
    \item it is the root of the whole game
    \item it and all of its successors form a self-contained part of the tree
    \item it is a terminal node
  \end{enumerate}
  \item In a game of perfect information, every information set ...
  \begin{enumerate}
    \item contains every node of the tree
    \item contains a single node
    \item is empty
    \item belongs to both players
  \end{enumerate}
  \item A strategy in an extensive form game specifies ...
  \begin{enumerate}
    \item only the first move
    \item the opponent's moves
    \item a single action for the whole game
    \item an action at every decision node of the player
  \end{enumerate}
  \item Subgame perfection rules out equilibria that rely on ...
  \begin{enumerate}
    \item cooperation
    \item non-credible threats off the path of play
    \item backward induction
    \item mixed strategies
  \end{enumerate}
  \item Backward induction at a decision node chooses the action that ...
  \begin{enumerate}
    \item maximises the other player's payoff
    \item minimises the total payoff
    \item maximises that player's payoff given the optimal play that follows
    \item is chosen at random
  \end{enumerate}
  \item If player 2 has a binary choice after each of player 1's two moves, player 2 has ...
  \begin{enumerate}
    \item two strategies
    \item eight strategies
    \item one strategy
    \item four contingent strategies
  \end{enumerate}
  \item Every subgame perfect equilibrium is ...
  \begin{enumerate}
    \item Pareto efficient
    \item unique
    \item a Nash equilibrium
    \item in mixed strategies
  \end{enumerate}
  \item Sequential rationality requires that a player's strategy is optimal ...
  \begin{enumerate}
    \item only at the first move
    \item only when the player is winning
    \item only on the equilibrium path
    \item at every information set, whether or not it is reached
  \end{enumerate}
  \item In games of perfect information, the equilibrium found by backward induction is ...
  \begin{enumerate}
    \item always mixed
    \item the worst outcome for both players
    \item never a Nash equilibrium
    \item subgame perfect
  \end{enumerate}
  \item An extensive form game has perfect information when ...
  \begin{enumerate}
    \item every information set contains a single node
    \item every player has the same payoff
    \item there are no chance moves
    \item players move simultaneously
  \end{enumerate}
  \item Backward induction proceeds by removing, at each decision node, actions that are ...
  \begin{enumerate}
    \item played first
    \item on the equilibrium path
    \item dominated for the player deciding there
    \item chosen by chance
  \end{enumerate}
  \item A subgame perfect equilibrium must induce a Nash equilibrium in ...
  \begin{enumerate}
    \item only the whole game
    \item every subgame, including those never reached in play
    \item only the final subgame
    \item only the subgames on the equilibrium path
  \end{enumerate}
  \item The requirement that a strategy be optimal at every information set where a player decides is called ...
  \begin{enumerate}
    \item additivity
    \item the folk theorem
    \item Pareto efficiency
    \item sequential rationality
  \end{enumerate}
  \item In every finite game of perfect information, backward induction produces ...
  \begin{enumerate}
    \item a non-credible threat
    \item no equilibrium in general
    \item a pure-strategy subgame perfect Nash equilibrium
    \item a mixed equilibrium only
  \end{enumerate}
\end{enumerate}
\end{document}
